\paragraph{The prescribed complex multiplication point.}
The period construction supplies the exact generic point needed for
the integral comparison; a separate choice of moduli component is
unnecessary.

\begin{proposition}\label{geo-cm-period-point-descent}
Let $X$ be a scheme of finite type over a number field $k$, carrying
an abelian family with coefficient action, its original symmetric
ample rigidified cubical line, and two prescribed finite level maps.
Assume that complex points of $X$ are separated by isomorphism classes
of these labelled objects, preserving both entire level maps.
Let $z\in X(\mathbf C)$ have a holomorphic period presentation
$A_z(\mathbf C)=\mathbf C^g/P\Lambda$, where $P$ identifies $\Lambda$
with a full period lattice, with the prescribed action of
an imaginary quadratic integer ring $\mathcal O_F$.
For its two embeddings $\iota_\pm$, set
\[
 \Lambda^\varepsilon=
 \{x\in\Lambda:P(ax)=\iota_\varepsilon(a)P(x)
                   \text{ for every }a\in\mathcal O_F\}.
\]
If $N\Lambda\subset\Lambda^++\Lambda^-$ for some integer $N>0$,
then the same point $z$ is defined over a number field containing $k$.
Its original line, coefficient action and finite level maps descend
with it. The field may contain any prescribed auxiliary number field.
\end{proposition}
\begin{proof}
The two rational pure subspaces have zero intersection, since the
embeddings of a nonreal element of $F$ differ. The finite-index
condition therefore gives pure integral vectors $x_1,\ldots,x_g$
forming an $F$-basis of $\Lambda\otimes\mathbf Q$.
Their $\mathcal O_F$-span has finite index in $\Lambda$, so the complex
span of the $P(x_j)$ contains a full period lattice and equals
$\mathbf C^g$.
The maps $w\mapsto\exp_{A_z}(wP(x_j))$ give an isogeny from a
power of the CM elliptic curve with period lattice
$\iota_+(\mathcal O_F)$ to the actual fibre $A_z$.
Use the conjugate coefficient action on a negative-pure factor.
The elliptic curve has algebraic $j$-invariant
\cite[Corollary~12.9]{cm-milne-mf}, and these holomorphic homomorphisms are algebraic
\cite[Theorems~2.8--2.9]{cm-milne-av}.

The finite kernel lies in algebraic torsion, so its quotient is defined
over $\overline{\mathbf Q}$. Homomorphism rigidity descends the
coefficient action and polarization. The exact symmetric-line descent
in Proposition~\ref{geo-cm-integral-product} retains the original line:
make an algebraic reference line symmetric by a square-root adjustment
in $\operatorname{Pic}^0$; its ratio with the specified line is then in
$\operatorname{Pic}^0[2]$. Normalize the comparison at zero, preserving
the cubical structure. The finite level maps have algebraic torsion
values. Thus the entire labelled object descends to
$\overline{\mathbf Q}$.

For every $\gamma\in\operatorname{Aut}(\mathbf C/\overline{\mathbf Q})$,
the objects at $z$ and $z^\gamma$ are labelled-isomorphic. Point
separation gives $z^\gamma=z$. On a $k$-affine open containing $z$,
every coordinate is therefore in $\overline{\mathbf Q}$. Indeed, a
transcendental coordinate $t$ can be included in a transcendence basis;
the automorphism $t\mapsto t+1$, fixing the other basis elements,
extends to $\mathbf C$. Finite type puts the coordinates in one number
field. Enlarge it to contain the desired auxiliary field.
Homomorphism rigidity and normalization at zero descend the product
maps and original-line comparison after a further finite extension.
The finite level maps agree by faithful base change.
The zero-dimensional case uses the empty elliptic
product and the same argument.
\end{proof}

For the fine moduli problem, uniqueness of the classifying map over
the reduced scheme $\operatorname{Spec}\mathbf C$ gives the stated
point separation. Let $\Lambda_0$ be the fixed lattice indexing the
level maps. If the first ideal is $I$ with $I\bar I=(M)$, then
$M\Lambda_0\subset I\Lambda_0$, and the second map is killed by
its specified positive integer. Thus both level maps are finite.
An existential compatible Tate system is not an additional labelled
tower. The proof compares two existing complex fibres, and therefore
requires no descent of such a tower.

The exact complex point is unchanged, so its chosen component and
original line are unchanged. Given the required proper arithmetic
model, its integral extension is supplied by
Proposition~\ref{geo-cm-integral-map-extension}.
The actual arithmetic moduli model and the pointwise period
presentation remain required; global surjectivity or algebraicity of
the period map is unnecessary for this point-descent argument.
